% !TEX root = thesis.tex
% ===============================================
% ===============================================
%	Abstract in german and english
% ===============================================
% ===============================================
\TUDoption{abstract}{multiple,section,notoc}
\renewcaptionname{english}{\abstractname}{Abstract}
\begin{abstract}[pagestyle=empty.tudheadings]%
Spiking Neural Networks (SNNs) offer a promising approach to energy-efficient and low-latency edge intelligence by processing information through discrete spike events. Field-Programmable Gate Arrays (FPGAs) are well suited to this model because they provide configurable parallelism and deterministic execution. However, most FPGA-based SNN accelerators use fixed synaptic weights, whereas online learning additionally requires neuron state, writable weight storage, and coordinated memory access. Previous neuromorphic processors demonstrate local learning but often rely on fixed, task-specific architectures, leaving a gap between specialized systems and reusable FPGA-generation frameworks.

This thesis addresses that gap by extending an accelerator generated by the open-source Spiker+ framework with a synthesizable Spike-Driven Synaptic Plasticity (SDSP) learning path. The Spiker+ hierarchy is retained while the postsynaptic states required by SDSP are exposed. The architecture combines local state evaluation, bounded updates, writable excitatory-weight memory, and an address-aligned read--modify--write mechanism that preserves the association between updated weights and their original addresses. Learning is selectable at runtime, and the same hardware retains an inference-only mode without issuing new weight writes.

The modified SNN and its learning and communication modules are packaged as a reusable Vivado IP core and deployed on a PYNQ-Z1. AXI4-Lite and AXI DMA connect the accelerator to a PYNQ runtime that prepares inputs, configures operation, transfers spike data, and retrieves classification results. Self-checking simulations, interface tests, and board-level experiments confirm correct supervision, persistent weight-dependent behaviour, and continuous operation across complete training epochs.

The implemented system meets timing at 100~MHz on the XC7Z020 and uses 5,056 lookup tables, 5,294 flip-flops, and four block RAMs. Including processor and transfer overhead, the complete application achieves approximately 66--68 images per second. In one continuous experiment, the accelerator completed 180,000 supervised transactions over three MNIST epochs. Strict accuracy on a fixed 1,000-image validation set increased from 5.9\% to 39.6\%. The single-layer 784--10 network provides a structural reference to TEDOP, while differences in models, precision, and measurement boundaries prevent a direct performance ranking. Overall, the results demonstrate persistent online adaptation and successful system-level integration into the reusable Spiker+ flow. They establish hardware feasibility, but not learning convergence, because only one parameter configuration and three epochs were evaluated. Longer training, parameter exploration, increased network capacity, and reduced processor-side overhead remain future work.
\end{abstract}
